%-----------------------------------------------------------------------------------------------
\makeatletter
\immediate\write18{datelog > \jobname.info} % site script for $(date '+%Y-%m-%d %Hh%Mm%Ss')
\makeatother
%-----------------------------------------------------------------------------------------------
%-----------------------------------------------------------------------------------------------
\usetheme{Copenhagen}
\usepackage{beamercolorthemeCNThermSci2}
\usefonttheme{serif}
%-----------------------------------------------------------------------------------------------

%-----------------------------------------------------------------------------------------------
%-----------------------------------------------------------------------------------------------
\usetheme{Copenhagen}
\usepackage{beamercolorthemeUTF2}
\usefonttheme{serif}
%-----------------------------------------------------------------------------------------------
\usepackage[utf8]{inputenc}
\usepackage[greek,french,english,brazil]{babel} % last becomes the active one
\usepackage{pslatex}
\usepackage{amssymb,amsmath}
\usepackage{soul}
\usepackage[squaren,Gray,cdot]{SIunits}
\usepackage[nice]{nicefrac}
\usepackage{tikz}
\usepackage{amscd}
\usepackage{stmaryrd}
\usepackage{scalerel}
\usepackage{xspace}
%-----------------------------------------------------------------------------------------------


%-----------------------------------------------------------------------------------------------
%-----------------------------------------------------------------------------------------------
% Mathematical
%-----------------------------------------------------------------------------------------------
\newcommand{\vet}[1]{\underline{{#1}}}
\newcommand{\mat}[1]{\underline{\underline{{#1}}}}
\newcommand{\cub}[1]{\underline{\underline{\underline{{#1}}}}}
\newcommand{\eqdef}{{\ensuremath\stackrel{\text{\tiny def}}{=}}}
%-----------------------------------------------------------------------------------------------
% Linguistic
%-----------------------------------------------------------------------------------------------
\newcommand{\GRtxt}[1]{\begin{otherlanguage}{greek}{{#1}}\end{otherlanguage}}
\newcommand{\FRtxt}[1]{\begin{otherlanguage}{french}{{#1}}\end{otherlanguage}}
%-----------------------------------------------------------------------------------------------
% Presentation
%-----------------------------------------------------------------------------------------------
\newcommand{\BkgImgH}[1]{% Places an image centered on the slide background filling the height
    \usebackgroundtemplate{\parbox{\paperwidth}{%
        \vspace*{1sp}\centering\includegraphics[height=\paperheight]{{#1}}
}}}
\newcommand{\BkgImgW}[1]{% Places an image centered on the slide background filling the width
    \usebackgroundtemplate{\parbox{\paperwidth}{%
        \vspace*{1sp}\centering\includegraphics[width=\paperwidth]{{#1}}
}}}
\newcommand{\ArtEndH}[3]{% Transitions to plain image (last) slide: #1:prefix #2,#3:extensions
    \BkgImgH{root/../art/#1.#2}
    \frame<handout:0>[plain]{%
        \transdissolve\vspace*{72mm}\color{white}\scriptsize\bf\input{root/../art/#1.#3}}
    \usebackgroundtemplate{\mbox{~}}
}
\newcommand{\ArtEndW}[3]{% Transitions to plain image (last) slide: #1:prefix #2,#3:extensions
    \BkgImgW{root/../art/#1.#2}
    \frame<handout:0>[plain]{%
        \transdissolve\vspace*{72mm}\color{white}\scriptsize\bf\input{root/../art/#1.#3}}
    \usebackgroundtemplate{\mbox{~}}
}
\newcommand{\ImgColW}[3]{% Inserts a full-width image in a column
    \includegraphics[width=\columnwidth]{root/../art/#1.#2}\\[-0.5\baselineskip]
    \parbox{\columnwidth}{\tiny\hfill\scalebox{0.85}{\input{root/../art/#1.#3}}}
}
\newcommand{\txtpic}[1]{%
    \fcolorbox{lightgray}{white!90!black}{{#1}} 
}
%-----------------------------------------------------------------------------------------------


%-----------------------------------------------------------------------------------------------
\title{B.01.01 -- Máquinas Hidráulicas de Fluxo}
\subtitle{Normas e Grandezas Básicas}
\author{Prof.~C.~Naaktgeboren, PhD}
\date{{\scriptsize\tt%
    \includegraphics[height=6.0mm]{cc/by-nc-nd-88x31.pdf}\\[\smallskipamount]
    https://github.com/CNThermSci/ApplThermSci\\
    Compiled on \input{\jobname.info}
}}
%-----------------------------------------------------------------------------------------------
\begin{document}
%-----------------------------------------------------------------------------------------------
\logo{%
    \parbox{158mm}{% There's a 1mm gap on each side of the 160mm x 90mm slide logo line
        \mode<beamer>{
            \includegraphics[height=6.0mm]{root/00-res/UTFPR/UTFPR-logo-D.pdf}\hfill%
            \includegraphics[height=9.0mm]{root/00-res/logo/CNThermSci-logo-A.pdf}%
        }
        \mode<handout>{
            \includegraphics[height=6.0mm]{root/00-res/UTFPR/UTFPR-logo-W.pdf}\hfill%
            \includegraphics[height=9.0mm]{root/00-res/logo/CNThermSci-logo-W.pdf}%
        }
    }
} % The (delineated, alpha), or washed-out logos
%-----------------------------------------------------------------------------------------------
\frame{\titlepage}
%-----------------------------------------------------------------------------------------------

%-----------------------------------------------------------------------------------------------
\frame{\tableofcontents}
%-----------------------------------------------------------------------------------------------

%-----------------------------------------------------------------------------------------------
\section{Normas em Máquinas de Fluxo}
%-----------------------------------------------------------------------------------------------

    % !j 96 -i8
    %-------------------------------------------------------------------------------------------
    \begin{frame}{Nomenclatura}\vspace*{-1em}

        \begin{itemize}
            \item<1-> Máquinas de fluxo são uma \alert{aplicação muito antiga} em fluidos;
            \item<1-> A \alert{nomenclatura} empregada é \alert{bastante heterogênea};
            \item<2-> Referências incluem \alert{acadêmicas} e \alert{industriais};
            \item<2-> \alert{Maiores fornecedores mundiais} convergem para o \alert{padrão IEC}.
            \item<3-> E também o material desta disciplina.
        \end{itemize}

    \end{frame}
    %-------------------------------------------------------------------------------------------

    % !j 96 -i8
    %-------------------------------------------------------------------------------------------
    \begin{frame}{Norma IEC 60193}\vspace*{-1em}

        IEC é o acrônimo da \textit{International Electrotechnical Commission\/}.
        \vspace*{\medskipamount}

        \begin{itemize}
            \item<1-> A \alert{\textbf{\texttt{IEC 60193:2019}}} cancela e revoga a
                $2^{\mathrm{a}}$ Ed.~de 1999;
            \item<1-> Aplica-se para \alert{modelos de laboratório} de máquinas de \alert{ação}
                e de \alert{reação};
            \item<1-> Aplica-se para \alert{turbinas hidráulicas}, \alert{bombas de
                armazenamento}, ou \alert{turbina-bombas};
            \item<2-> com \alert{potência unitária $> \unit{5}{\mega\watt}$}, ou
            \item<2-> com \alert{diâmetro $> \unit{3}{\meter}$};
            \item<3-> Esta norma objetiva \alert{definir termos e quantidades} empregados;
            \item<3-> além de estabelecer várias outras \alert{especificações},
                \alert{asserções} e \alert{garantias}...
        \end{itemize}

    \end{frame}
    %-------------------------------------------------------------------------------------------

    % !j 96 -i8
    %-------------------------------------------------------------------------------------------
    \begin{frame}{Outras Normas}\vspace*{-1em}

        \begin{itemize}
            \item<1-> \alert{IEC 60041} -- Field acceptance test to determine the hydraulic
                performance of hydraulic turbines, storage pumps and pump-turbines.

            \item<1-> \alert{IEC 60609} -- Cavitation pitting evaluation in hydraulic turbines,
                storage pumps and pump-turbines.

            \item<1-> \alert{IEC 60609-2} -- Cavitation pitting evaluation in hydraulic
                turbines, storage pumps and pump-turbines – Part 2: Evaluation in Pelton
                turbines.

            \item<1-> \alert{IEC 61364} -- Nomenclature of hydraulic machinery.

            \item<1-> \alert{VIM} -- International vocabulary of basic and general terms in
                metrology (ABNT ISO/IEC GUIA 99).
        \end{itemize}

    \end{frame}
    %-------------------------------------------------------------------------------------------

%-----------------------------------------------------------------------------------------------
\subsection{Definições -- IEC 60193}
%-----------------------------------------------------------------------------------------------

    % !j 96 -i8
    %-------------------------------------------------------------------------------------------
    \begin{frame}{Subscritos e Símbolos Pertinentes}\vspace*{-1em}

        \setlength{\tabcolsep}{2mm}
        \noindent\begin{longtable}{cp{110mm}}
            \alert{Símbolo} & \alert{Definição} \\
            \hline
            $1$ &
            Seção de referência de alta pressão \\
            $2$ &
            Seção de referência de baixa pressão \\
            $1'$ &
            Seção de medição de alta pressão \\
            $2'$ &
            Seção de medição de baixa pressão \\
            \hline
            max, min &
            Máximo ou mínimo valor, respectivamente \\
            $P$ &
            Referente ao \alert{protótipo}, em tamanho real \\
            $M$ &
            Referente ao \alert{modelo} em escala reduzida \\
            ref &
            Valores em condição de referência especificada \\
            amb &
            Valores referentes ao ambiente \\
            pl &
            Valores da planta \\
            $R$ &
            Referente à condição de \alert{disparo} (runaway) \\
            \hline
        \end{longtable}

    \end{frame}
    %-------------------------------------------------------------------------------------------

    % !j 96 -i8
    %-------------------------------------------------------------------------------------------
    \begin{frame}\vspace*{0em}
        \figfigH{fig/01-MachineRefs.pdf}
    \end{frame}
    %-------------------------------------------------------------------------------------------

    % !j 96 -i8
    %-------------------------------------------------------------------------------------------
    \begin{frame}{Termos Geométricos}\vspace*{-1em}

        \setlength{\tabcolsep}{2mm}
        \noindent\begin{longtable}{cp{110mm}}
            \alert{Símbolo} & \alert{Definição} \\
            \hline
            $A$ ($\meter^2$) &
            Área \\
            $a$ ($\meter$) &
            Abertura de palhetas (menor distância média entre palhetas adjacentes) \\
            $\alpha$ ($\degree$) &
            Ângulo de abertura de palhetas (valor médio à partir do fechamento) \\
            $\beta$ ($\degree$) &
            Ângulo de abertura de pá de rotor \\
            $D$ ($\meter$) &
            Diâmetro de referência (geralmente mínimo e não variável) \\
            $z$ ($\meter$) &
            Nível, ou quota (elevação em rel.~a uma ref.: nível do mar) \\
            \hline
        \end{longtable}

    \end{frame}
    %-------------------------------------------------------------------------------------------

    % !j 96 -i8
    %-------------------------------------------------------------------------------------------
    \begin{frame}\vspace*{-1em}
        \figfigH{fig/02-GuideVanes.pdf}
    \end{frame}
    %-------------------------------------------------------------------------------------------

    % !j 96 -i8
    %-------------------------------------------------------------------------------------------
    \begin{frame}{Quantidades e Propriedades Físicas}\vspace*{-1em}

        \setlength{\tabcolsep}{2mm}
        \noindent\begin{longtable}{cp{110mm}}
            \alert{Símbolo} & \alert{Definição} \\
            \hline
            $g$ ($\meter\per\second^2$) &
            Aceleração devido à gravidade \\
            $\Theta$ ($\kelvin$) &
            Temperatura termodinâmica \\
            $\theta$ ($\celsius$) &
            Temperatura em Celsius, $\theta = \Theta - 273,15$ \\
            $\rho$ ($\kilogram\per\meter^3$) &
            Densidade. Subscritos incluem: $_w$, $_a$ e $_{Hg}$, para água, ar e mercúrio \\
            \hline
        \end{longtable}

    \end{frame}
    %-------------------------------------------------------------------------------------------

%-----------------------------------------------------------------------------------------------
\section{Referências}
%-----------------------------------------------------------------------------------------------

    %------------------------------------------------------------------------------------------
    \begin{frame}{Referências}
        \begin{thebibliography}{Çengel, Y.~A., 2013}
            \bibitem[Potter, M. C., 2015]{2015-PotterMC-Cengage}
                Potter, M.~C., \textit{et al.\/}
                \newblock{%
                    {\em Mecânica dos Fluidos}.
                    \alert{Seção~12-1.}
                }
                \newblock{\footnotesize Cengage. São Paulo. ISBN 978-85-221-1568-6.}
            \bibitem[IEC60193]{IEC-60193:2019}
                International Electrotechnical Commission, \alert{IEC 60193}
                \newblock{%
                    {\em Hydraulic turbines, storage pumps and pump-turbines -- Model acceptance
                    tests}.
                }
                \newblock{\footnotesize International Standard. 2019.
                    \texttt{webstore.iec.ch/publication/60951}.}
        \end{thebibliography}
    \end{frame}
    %------------------------------------------------------------------------------------------

    % Finishes with stunning image, with credit
    \ArtEndWSE{pexels-mali-maeder-142497}{jpg}{txt}

%-----------------------------------------------------------------------------------------------
\end{document}
%-----------------------------------------------------------------------------------------------

